\chapter{Quantiles, Box Plots and Shape}
\label{ch:quantiles}

Chapter 2 summarised a data set with two numbers: a centre (mean or median) and a
spread (standard deviation or IQR). Two numbers cannot tell you that a service is fine
for most users and terrible for one in a hundred. This chapter adds the tools that can:
quantiles, which locate any position in the data; box plots and ECDFs, which draw
those positions; and a vocabulary for shape (skew, heavy tails, several bumps) with
the plots that reveal it.

Here is the running data set. A service handled twelve requests, with these
latencies in milliseconds, already sorted:
\[
  12,\ 14,\ 15,\ 16,\ 17,\ 18,\ 20,\ 22,\ 25,\ 31,\ 48,\ 120.
\]
The mean is $358/12 \approx 29.8$ ms and the median is $(18+20)/2 = 19$ ms. Only
three of the twelve requests were slower than the mean. One slow request (120 ms)
dragged the mean far above what a typical request experienced.

\begin{figure}[H]
\centering
\begin{tikzpicture}
\begin{axis}[
  width=12cm, height=3.6cm,
  xmin=0, xmax=130, ymin=0, ymax=2.2,
  axis lines=left, ytick=\empty, axis y line=none,
  xlabel={latency (ms)}, clip=false,
]
\addplot[only marks, mark=*, mark size=2.4pt, sourceblue] coordinates {
  (12,0.5) (14,0.5) (15,0.5) (16,0.5) (17,0.5) (18,0.5) (20,0.5) (22,0.5)
  (25,0.5) (31,0.5) (48,0.5) (120,0.5)};
\draw[targetgreen, thick] (axis cs:19,0.1) -- (axis cs:19,1.3);
\node[targetgreen, anchor=south east] at (axis cs:20,1.3) {\small median 19};
\draw[densityred, thick, dashed] (axis cs:29.83,0.1) -- (axis cs:29.83,1.75);
\node[densityred, anchor=south west] at (axis cs:29.83,1.75) {\small mean 29.8};
\end{axis}
\end{tikzpicture}
\caption{Twelve request latencies. Nine requests sit left of the mean. The single 120 ms
request pulls the mean right; the median stays with the bulk of the data.}
\label{fig:quantiles-dotplot}
\end{figure}

\section{Percentiles and quantiles}

The median splits the data in half. Nothing stops us from splitting at other points.

\begin{definition}[Quantile, percentile]
For $0 < p < 1$, a \textbf{$p$-quantile} of a data set is a number $q$ such that
roughly a fraction $p$ of the data lies at or below $q$ and a fraction $1-p$ lies
above it. The \textbf{$100p$-th percentile} is the same thing on a $0$--$100$ scale:
the $0.9$-quantile is the 90th percentile, written p90.
\end{definition}

Picture walking along the sorted data from left to right with a counter that tracks
the fraction of points you have passed. The $p$-quantile is where the counter reads
$p$. The median is the $0.5$-quantile. The \textbf{quartiles} are the $0.25$-, $0.5$-,
and $0.75$-quantiles, called $Q_1$, $Q_2$, $Q_3$: they cut the data into four
equal-sized groups.

\subsection{The trouble with ``roughly''}

With twelve points, what is p90? Ten percent of twelve is $1.2$ points. No data
value has exactly $1.2$ points above it. The word ``roughly'' in the definition is
hiding a real choice, and different software makes it differently.

The most common rule, used by default in numpy, R, Excel's \texttt{PERCENTILE.INC} and
pandas, places the sorted values at evenly spaced positions and draws straight lines
between them.

\begin{definition}[Linear-interpolation quantile]
Sort the data as $x_{(1)} \le x_{(2)} \le \cdots \le x_{(n)}$, where $x_{(i)}$ is
the $i$-th smallest value (the \textbf{$i$-th order statistic}). Put $x_{(i)}$ at
position $(i-1)/(n-1)$ on a $0$-to-$1$ scale, so the minimum sits at $0$ and the
maximum at $1$. Compute
\[
  h = (n-1)p + 1, \qquad k = \floor{h}, \qquad f = h - k .
\]
Then
\[
  Q(p) = x_{(k)} + f\,\bigl(x_{(k+1)} - x_{(k)}\bigr).
\]
\end{definition}

In words: $h$ is a fractional rank. If $h = 10.9$, go to the 10th value and then
$90\%$ of the way to the 11th.

\begin{example}[Quartiles and p90 of the twelve latencies]
\label{ex:quantiles-linear}
Here $n = 12$, so $h = 11p + 1$.
\begin{itemize}
  \item $Q_1$: $h = 11(0.25)+1 = 3.75$. Go to $x_{(3)} = 15$, then $75\%$ of the way
  to $x_{(4)} = 16$: $Q_1 = 15 + 0.75(1) = 15.75$.
  \item Median: $h = 6.5$. Halfway from $x_{(6)}=18$ to $x_{(7)}=20$: $19$.
  \item $Q_3$: $h = 9.25$. From $x_{(9)} = 25$, a quarter of the way to
  $x_{(10)} = 31$: $Q_3 = 25 + 0.25(6) = 26.5$.
  \item p90: $h = 10.9$. From $x_{(10)} = 31$, $90\%$ of the way to $x_{(11)} = 48$:
  $31 + 0.9(17) = 46.3$.
\end{itemize}
\end{example}

\begin{figure}[H]
\centering
\begin{tikzpicture}[x=0.95cm, y=1cm]
\draw[gray!60] (1,0) -- (12,0);
\foreach \i/\v in {1/12,2/14,3/15,4/16,5/17,6/18,7/20,8/22,9/25,10/31,11/48,12/120} {
  \fill[sourceblue] (\i,0) circle (2.4pt);
  \node[below, font=\small] at (\i,-0.12) {$\v$};
  \node[above, font=\scriptsize, gray] at (\i,0.12) {$\i$};
}
\node[left, font=\small, gray] at (0.6,0.32) {rank};
\node[left, font=\small] at (0.6,-0.4) {value};
\draw[-{Stealth}, quotientorange, thick] (3.75,1.35) -- (3.75,0.5);
\node[quotientorange, above, font=\small] at (3.75,1.35) {$h=3.75 \Rightarrow Q_1 = 15.75$};
\draw[-{Stealth}, targetgreen, thick] (6.5,0.95) -- (6.5,0.5);
\node[targetgreen, above, font=\small] at (6.5,0.95) {$h=6.5 \Rightarrow 19$};
\draw[-{Stealth}, quotientorange, thick] (9.25,0.95) -- (9.25,0.5);
\node[quotientorange, above, font=\small] at (9.25,0.95) {$h=9.25 \Rightarrow 26.5$};
\draw[-{Stealth}, densityred, thick] (10.9,1.75) -- (10.9,0.5);
\node[densityred, above, font=\small] at (10.9,1.75) {$h=10.9 \Rightarrow \text{p90} = 46.3$};
\end{tikzpicture}
\caption{Linear interpolation treats rank as a continuous scale. A fractional rank
$h$ lands between two data points, and the quantile is the matching point on the
straight line between their values.}
\label{fig:quantiles-interp}
\end{figure}

\subsection{Why numpy has nine methods}

The linear rule is one choice. Others put $x_{(i)}$ at position $i/(n+1)$, or at
$(i - 0.5)/n$, or refuse to interpolate and return an actual data value. Hyndman and
Fan (1996) catalogued nine such rules, and \texttt{numpy.quantile} exposes all of them
through its \texttt{method} argument. Table~\ref{tab:quantiles-methods} shows what
several of them return for our twelve latencies.

\begin{table}[H]
\centering
\begin{tabular}{lllcccc}
\toprule
numpy \texttt{method} & position of $x_{(i)}$ & R type & $Q_1$ & median & $Q_3$ & p90 \\
\midrule
\texttt{inverted\_cdf} & no interpolation & 1 & 15 & 18 & 25 & 48 \\
\texttt{lower} & round rank down & -- & 15 & 18 & 25 & 31 \\
\texttt{nearest} & round rank to nearest & -- & 16 & 20 & 25 & 48 \\
\texttt{hazen} & $(i-0.5)/n$ & 5 & 15.5 & 19 & 28 & 69.6 \\
\texttt{weibull} & $i/(n+1)$ & 6 & 15.25 & 19 & 29.5 & 98.4 \\
\texttt{linear} (default) & $(i-1)/(n-1)$ & 7 & 15.75 & 19 & 26.5 & 46.3 \\
\bottomrule
\end{tabular}
\caption{Six quantile rules applied to the same twelve numbers. They nearly agree in
the middle and disagree wildly at p90, where the data are sparse.}
\label{tab:quantiles-methods}
\end{table}

The \texttt{weibull} p90 is more than double the \texttt{linear} p90. Check it by hand:
with positions $i/(n+1)$ the fractional rank is $h = (n+1)p = 13(0.9) = 11.7$, so
$Q = x_{(11)} + 0.7\,(x_{(12)} - x_{(11)}) = 48 + 0.7(72) = 98.4$. The rule reaches
further into the gap between 48 and 120 because it reserves some probability for
values beyond the largest observation.

None of these rules is wrong. They answer slightly different questions (``which actual
data value?'', ``where, on average, does the $i$-th smallest of $n$ draws sit?'') and
they all converge to the same answer as $n$ grows. The gap matters when $n$ is small or $p$ is
extreme. A p99 computed from 50 requests rests on the one or two largest values, and
the method choice can swing it by a factor of two.

\begin{trapbox}
Two teams compute ``p90 latency'' from the same log and get different answers. Before
debugging the pipeline, check which quantile method each tool uses. Spreadsheets,
databases and monitoring systems do not all default to numpy's \texttt{linear}.
\end{trapbox}

\begin{lstlisting}[language=Python]
import numpy as np
x = np.array([12,14,15,16,17,18,20,22,25,31,48,120])
np.quantile(x, [0.25, 0.5, 0.75, 0.9])          # 15.75 19 26.5 46.3
np.quantile(x, 0.9, method="weibull")           # 98.4
np.percentile(x, 90)                            # same as quantile(x, 0.9)
\end{lstlisting}

\begin{recoverybox}
Without computing: as $n \to \infty$ with the same underlying population, do the
\texttt{linear} and \texttt{weibull} p90 values move closer together or further apart?
(Closer. The positions $(i-1)/(n-1)$ and $i/(n+1)$ differ by at most $1/(n+1)$.)
\end{recoverybox}

\section{The five-number summary and the box plot}

Five quantiles give a compact portrait of a data set:
\[
  \text{min},\quad Q_1,\quad \text{median},\quad Q_3,\quad \text{max}.
\]
For the latencies (linear method): $12,\ 15.75,\ 19,\ 26.5,\ 120$. The middle half of
the requests took between $15.75$ and $26.5$ ms. The \textbf{interquartile range} is the
width of that middle half:
\[
  \IQR = Q_3 - Q_1 = 26.5 - 15.75 = 10.75 \text{ ms}.
\]
The five numbers already tell a story about shape. The distance from the median up to
$Q_3$ (7.5 ms) is more than double the distance down to $Q_1$ (3.25 ms), and the max is
93.5 ms above $Q_3$ while the min is only 3.75 ms below $Q_1$. The data stretch to
the right.

\subsection{Drawing the box plot}

A \textbf{box plot} turns the five-number summary into a picture, with one extra rule
for flagging unusual points.

\begin{enumerate}
  \item Draw a box from $Q_1$ to $Q_3$. Its length is the IQR.
  \item Draw a line across the box at the median.
  \item Compute the \textbf{fences}, invisible boundaries $1.5$ IQRs beyond the box:
  \[
    \text{lower fence} = Q_1 - 1.5\,\IQR, \qquad \text{upper fence} = Q_3 + 1.5\,\IQR.
  \]
  \item Extend a \textbf{whisker} from each end of the box to the most extreme data
  point that is still inside the fences.
  \item Plot every point beyond a fence individually. These are flagged as
  \textbf{outliers}: points far enough from the middle half to deserve a look.
\end{enumerate}

For the latencies: $1.5 \times 10.75 = 16.125$, so the lower fence is
$15.75 - 16.125 = -0.375$ and the upper fence is $26.5 + 16.125 = 42.625$. Nothing
falls below the lower fence, so the lower whisker reaches the minimum, 12. The largest
value inside the upper fence is 31, so the upper whisker stops there. The values 48
and 120 are plotted as outliers.

\begin{figure}[H]
\centering
\begin{tikzpicture}
\begin{axis}[
  width=12.5cm, height=5.6cm,
  xmin=-5, xmax=130, ymin=0.2, ymax=2.4,
  axis lines=left, axis y line=none,
  xlabel={latency (ms)}, clip=false,
  boxplot/draw direction=x,
]
\addplot+[boxplot prepared={lower whisker=12, lower quartile=15.75, median=19,
  upper quartile=26.5, upper whisker=31, draw position=1, box extend=0.45},
  draw=sourceblue, fill=histfill!50, thick, mark=none] coordinates {};
\addplot[only marks, mark=o, mark size=2.6pt, thick, trapred] coordinates {(48,1) (120,1)};
\draw[trapred, dashed] (axis cs:42.625,0.95) -- (axis cs:42.625,1.75);
\draw[trapred, dashed] (axis cs:-0.375,0.95) -- (axis cs:-0.375,1.75);
\node[trapred, font=\scriptsize, anchor=south west] at (axis cs:40,1.75) {upper fence 42.6};
\node[trapred, font=\scriptsize, anchor=south] at (axis cs:-0.375,1.75) {lower fence $-0.4$};
\draw[gray] (axis cs:15.75,1.25) -- (axis cs:15.75,1.55);
\draw[gray] (axis cs:19,1.25) -- (axis cs:19,1.95);
\draw[gray] (axis cs:26.5,1.25) -- (axis cs:26.5,1.55);
\node[font=\scriptsize, anchor=south east] at (axis cs:16.5,1.55) {$Q_1$};
\node[font=\scriptsize, anchor=south] at (axis cs:19,1.95) {median};
\node[font=\scriptsize, anchor=south west] at (axis cs:25.75,1.55) {$Q_3$};
\draw[decorate, decoration={brace, amplitude=4pt, mirror}] (axis cs:15.75,0.7) -- (axis cs:26.5,0.7)
  node[midway, below=4pt, font=\scriptsize] {IQR $=10.75$};
\node[font=\scriptsize, anchor=west] at (axis cs:32,0.86) {whisker end 31};
\node[trapred, font=\scriptsize, anchor=north] at (axis cs:120,0.88) {outlier};
\node[trapred, font=\scriptsize, anchor=south west] at (axis cs:46,1.08) {outlier};
\end{axis}
\end{tikzpicture}
\caption{Anatomy of a box plot for the twelve latencies. The whiskers stop at the last
data points inside the fences (12 and 31). The two slow requests are drawn
individually because they lie beyond the upper fence.}
\label{fig:quantiles-box-anatomy}
\end{figure}

\begin{whybox}
Why $1.5$? John Tukey chose it so that clean bell-shaped data would almost never
trigger the rule. For a Normal distribution with standard deviation $\sigma$, the
quartiles sit $0.674\sigma$ from the centre, so $\IQR = 1.349\sigma$ and the upper fence
is at $0.674\sigma + 1.5(1.349\sigma) \approx 2.70\sigma$. The chance of a Normal value
beyond either fence is about $0.7\%$. A flagged point in bell-shaped data is therefore
genuinely rare: roughly 7 per 1000 points.
\end{whybox}

\begin{trapbox}
``Outlier'' in a box plot means ``beyond the fences'', not ``error'' or ``delete me''. For
skewed data the rule fires constantly. Exponential data (a common model for waiting
times) have $Q_1 = \ln(4/3)$ and $Q_3 = \ln 4$ in units of the mean, so the upper
fence is $\ln 4 + 1.5\ln 3 \approx 3.03$ and about $4.8\%$ of all points lie beyond it.
A service with 10{,}000 requests would show roughly 480 ``outliers'' that are simply
its normal tail.
\end{trapbox}

Two more conventions you will meet. Some textbooks compute $Q_1$ and $Q_3$ as the
medians of the lower and upper halves of the data (Tukey's \emph{hinges}). For our
twelve points the halves are $\{12,\dots,18\}$ and $\{20,\dots,120\}$, giving $15.5$
and $28$: slightly different from $15.75$ and $26.5$, for the reason in the previous
section. And some plotting libraries draw whiskers to the min and max with no
outlier rule. Read the documentation before reading the whiskers.

\subsection{Comparing groups side by side}

The real power of the box plot is comparison. Put several box plots on one axis and
your eye compares medians, spreads and tails at once. Below, 60 requests were sampled
from each of three regions.

\begin{figure}[H]
\centering
\begin{tikzpicture}
\begin{axis}[
  width=11cm, height=6.4cm,
  ymin=0, ymax=165, xmin=0.4, xmax=3.6,
  xtick={1,2,3}, xticklabels={us-east, eu-west, ap-south},
  ylabel={latency (ms)}, axis lines=left,
  boxplot/draw direction=y,
]
\addplot+[boxplot prepared={lower whisker=7.4, lower quartile=17.03, median=19.2,
  upper quartile=24.5, upper whisker=34.4, draw position=1, box extend=0.4},
  draw=sourceblue, fill=histfill!50, thick, mark=none] coordinates {};
\addplot+[boxplot prepared={lower whisker=14.1, lower quartile=22.0, median=28.25,
  upper quartile=35.1, upper whisker=46.4, draw position=2, box extend=0.4},
  draw=targetgreen, fill=targetgreen!20, thick, mark=none] coordinates {};
\addplot+[boxplot prepared={lower whisker=10.0, lower quartile=26.27, median=37.45,
  upper quartile=56.62, upper whisker=100.6, draw position=3, box extend=0.4},
  draw=quotientorange, fill=quotientorange!20, thick, mark=none] coordinates {};
\addplot[only marks, mark=o, thick, sourceblue] coordinates {(1,39.4) (1,40.9) (1,64.0)};
\addplot[only marks, mark=o, thick, targetgreen] coordinates {(2,60.2)};
\addplot[only marks, mark=o, thick, quotientorange] coordinates {(3,119.6) (3,128.0) (3,153.4)};
\end{axis}
\end{tikzpicture}
\caption{Latency by region, 60 requests each. ap-south has the highest median and by
far the widest box: its typical request is slower and much less predictable. Every
flagged point is on the high side, a sign of right skew.}
\label{fig:quantiles-box-groups}
\end{figure}

Read it from the boxes outward. Medians: about $19$, $28$ and $37$ ms. IQRs: $7.5$,
$13.1$ and $30.4$ ms. The ap-south box is four times as tall as us-east's, which matters
more to users than the 18 ms difference in medians. Each region has a few flagged
points, all on the high side. A table of means ($21.4$, $29.1$, $45.6$) would have
hidden almost all of this.

\section{The empirical CDF}

A box plot keeps five numbers. The \textbf{empirical cumulative distribution
function} (ECDF) keeps all of them.

\begin{definition}[Empirical CDF]
For data $x_1, \dots, x_n$, the ECDF is the function
\[
  \hat F(t) = \frac{\#\{i : x_i \le t\}}{n},
\]
the fraction of data points at or below $t$.
\end{definition}

Picture sliding a vertical line from left to right across a dot plot. Each time it
passes a data point, the ECDF jumps up by $1/n$. Between data points it is flat. It
starts at $0$ far to the left and ends at $1$ far to the right. With our twelve
latencies each jump is $1/12 \approx 0.083$. For example $\hat F(20) = 7/12 \approx 0.58$
because seven requests took 20 ms or less.

\begin{figure}[H]
\centering
\begin{tikzpicture}
\begin{axis}[
  width=12cm, height=6.2cm,
  xmin=0, xmax=130, ymin=0, ymax=1.08,
  axis lines=left, xlabel={latency $t$ (ms)}, ylabel={$\hat F(t)$},
  ytick={0,0.25,0.5,0.75,0.9,1}, clip=false,
]
\addplot[const plot, sourceblue, thick] coordinates {
  (0,0) (12,0.0833) (14,0.1667) (15,0.25) (16,0.3333) (17,0.4167) (18,0.5)
  (20,0.5833) (22,0.6667) (25,0.75) (31,0.8333) (48,0.9167) (120,1) (130,1)};
\addplot[only marks, mark=*, mark size=1.6pt, sourceblue] coordinates {
  (12,0.0833) (14,0.1667) (15,0.25) (16,0.3333) (17,0.4167) (18,0.5)
  (20,0.5833) (22,0.6667) (25,0.75) (31,0.8333) (48,0.9167) (120,1)};
\draw[densityred, dashed] (axis cs:0,0.9) -- (axis cs:48,0.9) -- (axis cs:48,0);
\node[densityred, anchor=south west, font=\small] at (axis cs:50,0.78) {$p=0.9$ first reached at $t=48$};
\draw[targetgreen, dashed] (axis cs:0,0.5) -- (axis cs:18,0.5) -- (axis cs:18,0);
\node[targetgreen, anchor=west, font=\small] at (axis cs:50,0.45) {$p=0.5$ first reached at $t=18$};
\end{axis}
\end{tikzpicture}
\caption{ECDF of the twelve latencies. Read a quantile by going across from $p$ on the
vertical axis until you hit the staircase, then down. The long flat step from 48 to 120
is the gap in the data that made the quantile methods disagree.}
\label{fig:quantiles-ecdf}
\end{figure}

To read a quantile off the ECDF, go horizontally from height $p$ until you meet the
staircase, then drop to the horizontal axis. At $p = 0.9$ the staircase first reaches
$0.9$ at $t = 48$ (where it jumps from $0.83$ to $0.92$). That is the
\texttt{inverted\_cdf} answer in Table~\ref{tab:quantiles-methods}: the ``no
interpolation'' rule is exactly ``invert the ECDF''. The interpolating rules replace the
staircase with a sloped line through it, which is why they produce values between data
points.

The ECDF has three advantages over a histogram. It needs no bin width. Every data
point is visible. And quantiles, medians and ``fraction below a threshold'' are all
read directly. A service-level objective (SLO), a promise such as ``most requests
finish under 30 ms'', becomes the question ``what fraction finished under 30 ms?'', and
the answer is $\hat F(30) = 9/12 = 0.75$, one glance.

\begin{recoverybox}
Two ECDFs are drawn on one plot. Curve A lies everywhere to the left of curve B. Which
data set has the smaller median? (A. At height 0.5, A is reached first.)
\end{recoverybox}

\section{The quantile function: the CDF read backwards}

The same picture works for a whole population, not just a sample. Think of a random
variable $X$ as the latency of the next request, before it happens. Its
\textbf{cumulative distribution function} (CDF) is $F(x) = \P(X \le x)$: the
probability that the request finishes in $x$ ms or less. It is the population version of
the ECDF: a curve rising from $0$ to $1$, smooth instead of a staircase. The ECDF of a
large sample hugs it closely. (\emph{Probability From the Ground Up} develops CDFs in
full.)

\begin{definition}[Quantile function]
The \textbf{quantile function} of $X$ is
\[
  Q(p) = \min\{x : F(x) \ge p\}, \qquad 0 < p < 1 ,
\]
the smallest value at which the CDF reaches $p$. When $F$ is continuous and strictly
increasing, $Q$ is the inverse function $F^{-1}$.
\end{definition}

The CDF answers ``given a value, what fraction lies below it?'' The quantile function
answers the reverse: ``given a fraction, what value has that fraction below it?'' Their
graphs are mirror images across the diagonal $y = x$.

\begin{example}[Exponential latencies]
Suppose latency is exponential with mean $20$ ms, so $F(x) = 1 - e^{-x/20}$ for $x \ge 0$.
Solve $F(x) = p$ for $x$:
\[
  1 - e^{-x/20} = p \;\Longrightarrow\; Q(p) = -20\ln(1-p).
\]
The median is $Q(0.5) = 20\ln 2 \approx 13.9$ ms, well below the mean of 20. The p99 is
$Q(0.99) = 20\ln 100 \approx 92.1$ ms, and the p99.9 is $20\ln 1000 \approx 138.2$ ms.
Each extra ``nine'' adds the same $20\ln 10 \approx 46$ ms.
\end{example}

\begin{figure}[H]
\centering
\begin{tikzpicture}
\begin{axis}[
  name=cdf, width=6.3cm, height=5.6cm,
  xmin=0, xmax=110, ymin=0, ymax=1.05,
  axis lines=left, xlabel={$x$ (ms)}, ylabel={$F(x)$},
  title={CDF: value $\to$ fraction}, title style={font=\small},
  domain=0:110, samples=80,
]
\addplot[distblue, thick] {1-exp(-x/20)};
\draw[densityred, -{Stealth}] (axis cs:0,0.9) -- (axis cs:46.05,0.9);
\draw[densityred, -{Stealth}] (axis cs:46.05,0.9) -- (axis cs:46.05,0);
\node[densityred, anchor=south west, font=\scriptsize] at (axis cs:48,0.02) {46.1};
\node[densityred, anchor=south, font=\scriptsize] at (axis cs:20,0.9) {$p=0.9$};
\end{axis}
\begin{axis}[
  at={(cdf.east)}, anchor=west, xshift=1.4cm,
  width=6.3cm, height=5.6cm,
  xmin=0, xmax=1.02, ymin=0, ymax=110,
  axis lines=left, xlabel={$p$}, ylabel={$Q(p)$ (ms)},
  title={Quantile function: fraction $\to$ value}, title style={font=\small},
  domain=0:0.995, samples=100,
]
\addplot[quotientorange, thick] {-20*ln(1-x)};
\draw[densityred, dashed] (axis cs:0.9,0) -- (axis cs:0.9,46.05) -- (axis cs:0,46.05);
\draw[densityred, dashed] (axis cs:0.99,0) -- (axis cs:0.99,92.1) -- (axis cs:0,92.1);
\node[densityred, anchor=south east, font=\scriptsize] at (axis cs:0.88,46.05) {p90 $=46.1$};
\node[densityred, anchor=south east, font=\scriptsize] at (axis cs:0.88,92.1) {p99 $=92.1$};
\end{axis}
\end{tikzpicture}
\caption{The exponential CDF with mean 20 ms (left) and its quantile function (right).
Same curve, axes swapped. The quantile function shoots up as $p \to 1$: the tail
quantiles are where the large values live.}
\label{fig:quantiles-qf}
\end{figure}

Quantiles have one property the mean lacks: they pass straight through any increasing
transformation. If $g$ is increasing (such as $\log$ or $\sqrt{\phantom{x}}$), then
\[
  Q_{g(X)}(p) = g\bigl(Q_X(p)\bigr),
\]
because $g$ preserves the order of the data, and quantiles only depend on order. The
median of log-latency is the log of the median latency. The mean has no such rule:
the mean of the logs is not the log of the mean.

\begin{connectionbox}
The quantile function is also how computers draw random numbers from a given
distribution. If $U$ is uniform on $(0,1)$, then $Q(U)$ has CDF $F$. This is
\textbf{inverse transform sampling}: \texttt{-20*np.log(1-np.random.rand())} draws an
exponential latency with mean 20. The proof is one line, $\P(Q(U) \le x) =
\P(U \le F(x)) = F(x)$; see \emph{Probability From the Ground Up} for details. The same
idea powers quantile regression and the ``pinball loss'' used to train models that
predict p90 instead of the mean.
\end{connectionbox}

\section{Latency dashboards: p50, p90, p99}

Production monitoring rarely shows a mean latency. It shows a few percentiles, usually
p50, p90 and p99, sometimes p99.9. Here is why.

Take 200 requests from a typical web service. Figure~\ref{fig:quantiles-lat-hist}
shows their histogram with four summary numbers marked.

\begin{figure}[H]
\centering
\begin{tikzpicture}
\begin{axis}[
  width=12.5cm, height=6cm,
  xmin=0, xmax=200, ymin=0, ymax=72,
  axis lines=left, xlabel={latency (ms)}, ylabel={requests},
  clip=false, axis on top,
]
\addplot[ybar interval, fill=histfill, draw=white] coordinates {
  (0,13) (10,60) (20,50) (30,25) (40,20) (50,11) (60,3) (70,7) (80,4) (90,2)
  (100,0) (110,1) (120,3) (130,0) (140,0) (150,0) (160,0) (170,0) (180,0) (190,1) (200,0)};
\draw[targetgreen, very thick] (axis cs:24.75,0) -- (axis cs:24.75,66);
\node[targetgreen, anchor=south, font=\small] at (axis cs:20,66) {p50 $=24.8$};
\draw[densityred, very thick, dashed] (axis cs:32.45,0) -- (axis cs:32.45,58);
\node[densityred, anchor=west, font=\small] at (axis cs:33,58) {mean $=32.5$};
\draw[quotientorange, very thick] (axis cs:60.87,0) -- (axis cs:60.87,40);
\node[quotientorange, anchor=west, font=\small] at (axis cs:61.5,40) {p90 $=60.9$};
\draw[structurepurple, very thick] (axis cs:124.73,0) -- (axis cs:124.73,25);
\node[structurepurple, anchor=west, font=\small] at (axis cs:125.5,25) {p99 $=124.7$};
\end{axis}
\end{tikzpicture}
\caption{200 request latencies. The mean (32.5 ms) is not a latency anyone typically
gets: about two thirds of requests are faster. The p99 (124.7 ms) is five times the
median and is invisible in both the mean and the median.}
\label{fig:quantiles-lat-hist}
\end{figure}

Each percentile answers a concrete question. p50: what does a typical request feel
like? p90: what do the unlucky 1-in-10 get? p99: what is the experience of the slowest
1\%? The mean answers none of these. It mixes the typical and the extreme into one
number that describes neither.

\subsection{Why averages lie about tails}

Figure~\ref{fig:quantiles-dashboard} shows a 30-minute dashboard, one point per minute,
each summarising 600 requests. Between minutes 12 and 17, 2\% of requests hit a slow
code path that takes about 400 ms.

\begin{figure}[H]
\centering
\begin{tikzpicture}
\begin{axis}[
  width=12.5cm, height=6.2cm,
  xmin=0, xmax=29, ymin=10, ymax=1000, ymode=log,
  axis lines=left, xlabel={minute}, ylabel={latency (ms), log scale},
  legend style={font=\small}, legend pos=north west,
  log ticks with fixed point,
]
\fill[trapred!8] (axis cs:11.5,10) rectangle (axis cs:17.5,1000);
\addplot[structurepurple, thick, mark=*, mark size=1.2pt] coordinates {
  (0,74.0) (1,67.3) (2,78.2) (3,77.6) (4,79.5) (5,78.2) (6,74.7) (7,81.1) (8,69.3) (9,79.1)
  (10,82.2) (11,76.3) (12,391.5) (13,399.0) (14,412.3) (15,391.8) (16,401.0) (17,413.5)
  (18,87.2) (19,80.4) (20,65.0) (21,80.3) (22,71.6) (23,79.5) (24,77.6) (25,77.7)
  (26,71.3) (27,80.9) (28,78.7) (29,80.5)};
\addlegendentry{p99}
\addplot[densityred, thick, dashed, mark=square*, mark size=1.1pt] coordinates {
  (0,28.8) (1,28.6) (2,28.9) (3,27.9) (4,28.4) (5,28.1) (6,27.9) (7,28.7) (8,28.2) (9,28.6)
  (10,28.6) (11,27.5) (12,35.9) (13,36.0) (14,35.1) (15,35.7) (16,35.9) (17,36.9)
  (18,28.8) (19,28.5) (20,27.8) (21,28.7) (22,28.0) (23,28.3) (24,28.4) (25,27.8)
  (26,27.4) (27,28.5) (28,27.8) (29,29.7)};
\addlegendentry{mean}
\addplot[targetgreen, thick, mark=triangle*, mark size=1.3pt] coordinates {
  (0,25.6) (1,26.0) (2,25.5) (3,24.2) (4,25.1) (5,24.8) (6,24.6) (7,24.8) (8,25.4) (9,24.6)
  (10,25.0) (11,24.4) (12,26.3) (13,25.2) (14,25.1) (15,24.7) (16,25.4) (17,25.5)
  (18,25.5) (19,25.1) (20,24.9) (21,25.1) (22,25.4) (23,25.0) (24,24.5) (25,24.6)
  (26,24.6) (27,24.8) (28,25.0) (29,26.1)};
\addlegendentry{p50}
\end{axis}
\end{tikzpicture}
\caption{A slow path hits 2\% of requests during the shaded minutes. The median does
not move. The mean rises from about 28 to 36 ms, a change easy to dismiss as noise. The
p99 jumps from about 77 to 400 ms, unmistakable.}
\label{fig:quantiles-dashboard}
\end{figure}

The arithmetic explains the picture. Moving 2\% of requests from about 28 ms to about
400 ms raises the mean by roughly $0.02 \times (400 - 28) \approx 7.4$ ms. The median
cannot move at all, since 98\% of requests are unchanged. The p99 sits right inside the
slow 2\%, so it jumps to about 400 ms.

Tails matter more than their size suggests because users make many requests. If loading
one page fans out to 100 backend calls and the page waits for all of them, the chance
that at least one call is slower than the backend's p99 is
\[
  1 - 0.99^{100} \approx 0.63 .
\]
Most page loads see a p99 event. The backend's tail becomes the user's typical
experience.

\begin{trapbox}
A green mean is not evidence that users are happy. A mean can hide a slow 1\% entirely
and a slow 5\% largely. If you only have room for one latency number, use a high
percentile, and always pair it with the request count it was computed from.
\end{trapbox}

\subsection{Percentiles do not average}

A fleet of servers each reports its own p99. It is tempting to average them. Try the
same move on medians with tiny numbers. Server A served requests of $1, 2, 3$ ms
(median 2). Server B served $10, 11, 12, 13, 14$ ms (median 12). The average of the
medians is 7. But the fleet's eight requests $1, 2, 3, 10, 11, 12, 13, 14$ have median
$(10 + 11)/2 = 10.5$. The answer depends on how many requests each server handled and
how their distributions overlap, and the two medians alone do not carry that
information.

At scale the error is just as real. Suppose server A handles 9{,}000 fast requests
(p99 $= 52.0$ ms) and server B handles 1{,}000 slow ones (p99 $= 286.8$ ms). The average
of the two p99s is $169.4$ ms. The true fleet p99, computed from all 10{,}000 requests,
is $218.1$ ms. And $3.5\%$ of all requests exceed $169.4$ ms, so the ``averaged p99''
is really about a p96.5.

\begin{figure}[H]
\centering
\begin{tikzpicture}
\begin{axis}[
  width=12.5cm, height=3.8cm,
  xmin=0, xmax=320, ymin=0, ymax=3.6,
  axis lines=left, axis y line=none, xlabel={latency (ms)}, clip=false,
]
\draw[sourceblue, line width=5pt] (axis cs:8,2.6) -- (axis cs:52,2.6);
\node[sourceblue, anchor=east, font=\small] at (axis cs:7,2.6) {A};
\draw[quotientorange, line width=5pt] (axis cs:76,1.6) -- (axis cs:286.8,1.6);
\node[quotientorange, anchor=east, font=\small] at (axis cs:74,1.6) {B};
\draw[sourceblue, thick] (axis cs:52,2.2) -- (axis cs:52,3.0);
\node[sourceblue, anchor=south, font=\scriptsize] at (axis cs:52,3.0) {A p99 52};
\draw[quotientorange, thick] (axis cs:286.8,1.2) -- (axis cs:286.8,2.0);
\node[quotientorange, anchor=south, font=\scriptsize] at (axis cs:286.8,2.0) {B p99 287};
\draw[trapred, thick, dashed] (axis cs:169.4,0) -- (axis cs:169.4,3.2);
\node[trapred, anchor=south, font=\scriptsize] at (axis cs:150,3.2) {average of p99s 169};
\draw[targetgreen, very thick] (axis cs:218.1,0) -- (axis cs:218.1,3.2);
\node[targetgreen, anchor=south west, font=\scriptsize] at (axis cs:200,3.2) {true fleet p99 218};
\end{axis}
\end{tikzpicture}
\caption{Bars show each server's middle 98\% (p1 to p99), roughly. The fleet's slowest
1\% of requests (100 requests) all come from B, so the fleet p99 is B's p90, about 218 ms,
not the midpoint of the two p99s.}
\label{fig:quantiles-fleet}
\end{figure}

The fix is to merge data, not summaries. Either combine the raw latencies, or have every
server report a \emph{histogram} (counts in fixed latency buckets) or a mergeable sketch
such as a t-digest or HDR histogram. Bucket counts add across servers and across time;
the combined percentile is then read off the merged counts. The same rule applies over
time: the p99 for an hour is not the average of sixty per-minute p99s.

\section{Shape: skewness and kurtosis}

The latency histograms all lean the same way: a big pile on the left and a long tail
trailing right. That lean has a name and a number.

\begin{figure}[H]
\centering
\begin{tikzpicture}
\begin{axis}[
  name=left, width=4.7cm, height=4cm, axis lines=left, ytick=\empty,
  xmin=0, xmax=1, ymin=0, ymax=2.2, xtick=\empty,
  title={left-skewed}, title style={font=\small}, domain=0:1, samples=80,
]
\addplot[fill=histfill!50, draw=distblue, thick] {20*x^4*(1-x)} \closedcycle;
\draw[densityred, dashed] (axis cs:0.714,0) -- (axis cs:0.714,2.0);
\draw[targetgreen] (axis cs:0.736,0) -- (axis cs:0.736,2.0);
\end{axis}
\begin{axis}[
  name=mid, at={(left.east)}, anchor=west, xshift=0.25cm,
  width=4.7cm, height=4cm, axis lines=left, ytick=\empty,
  xmin=-3, xmax=3, ymin=0, ymax=0.48, xtick=\empty,
  title={symmetric}, title style={font=\small}, domain=-3:3, samples=80,
]
\addplot[fill=histfill!50, draw=distblue, thick] {0.39894*exp(-x^2/2)} \closedcycle;
\draw[densityred, dashed] (axis cs:0,0) -- (axis cs:0,0.44);
\end{axis}
\begin{axis}[
  name=right, at={(mid.east)}, anchor=west, xshift=0.25cm,
  width=4.7cm, height=4cm, axis lines=left, ytick=\empty,
  xmin=0, xmax=1, ymin=0, ymax=2.2, xtick=\empty,
  title={right-skewed}, title style={font=\small}, domain=0:1, samples=80,
]
\addplot[fill=histfill!50, draw=distblue, thick] {20*x*(1-x)^4} \closedcycle;
\draw[densityred, dashed] (axis cs:0.286,0) -- (axis cs:0.286,2.0);
\draw[targetgreen] (axis cs:0.264,0) -- (axis cs:0.264,2.0);
\end{axis}
\end{tikzpicture}
\caption{Skew is named after the direction of the long tail. Dashed red: mean; solid
green: median. In a right-skewed distribution the tail pulls the mean to the right of
the median; in a left-skewed one, to the left.}
\label{fig:quantiles-skew-shapes}
\end{figure}

A distribution is \textbf{right-skewed} (positively skewed) when its long tail points
right, toward large values. Latencies, incomes, house prices, file sizes and city
populations are right-skewed: they are bounded below (by zero) but not above. Exam
scores on an easy test are \textbf{left-skewed}: most students near the top, a tail of
low scores.

\begin{definition}[Sample skewness]
With $\xbar$ the mean, let $m_k = \frac1n\sum_i (x_i - \xbar)^k$ be the average $k$-th
power of the deviations. The sample \textbf{skewness} is
\[
  g_1 = \frac{m_3}{m_2^{3/2}} .
\]
\end{definition}

Why the cube? Cubing keeps the sign of each deviation, so points above the mean push
$g_1$ up and points below push it down. Cubing also exaggerates big deviations, so a
long tail on one side wins. Dividing by $m_2^{3/2}$ (the cube of a standard deviation)
makes the number unit-free: skewness in milliseconds and in seconds is the same.
Symmetric data give $g_1 \approx 0$.

For the twelve latencies, $g_1 \approx 2.52$: strongly right-skewed. The single value
120 ms supplies all of it. Its cubed deviation, $(120 - 29.83)^3 \approx 733{,}000$, is
larger than the whole sum of cubes ($\approx 718{,}000$). The nine fast requests
contribute negative cubes totalling about $-20{,}900$, more than cancelling the
$+6{,}000$ from 31 and 48.

\begin{intuitionbox}
Skewness is a measure of tail imbalance, and it is fragile. One extreme value can
dominate it, as it did here. For small samples treat $g_1$ as a description of the data
in hand, not a stable property of the population. A box plot (whisker lengths, median
position in the box) shows the same imbalance more robustly.
\end{intuitionbox}

\subsection{Kurtosis: how heavy are the tails?}

Skewness compares the two tails. \textbf{Kurtosis} measures how heavy both tails are
compared with a Normal distribution of the same standard deviation:
\[
  \text{excess kurtosis} = \frac{m_4}{m_2^2} - 3 .
\]
The fourth power amplifies large deviations even more than the cube, and ignores their
sign. The ``$-3$'' is there because a Normal distribution has $m_4/m_2^2 = 3$, so excess
kurtosis is $0$ for a Normal, positive for heavier tails, negative for lighter ones.

\begin{figure}[H]
\centering
\begin{tikzpicture}
\begin{axis}[
  width=12cm, height=4.8cm,
  xmin=-5, xmax=5, ymin=0, ymax=0.75,
  axis lines=left, xlabel={value (both have SD 1)}, ylabel={density},
  domain=-5:5, samples=161,
  legend style={font=\small, at={(0.5,1.03)}, anchor=south}, legend columns=2,
]
\addplot[distblue, thick] {0.39894*exp(-x^2/2)};
\addlegendentry{Normal, excess kurtosis 0}
\addplot[densityred, thick, dashed] {0.70711*exp(-1.41421*abs(x))};
\addlegendentry{Laplace, excess kurtosis 3}
\draw[gray, -{Stealth}] (axis cs:3.6,0.12) -- (axis cs:3.2,0.025);
\node[gray, anchor=south, font=\scriptsize] at (axis cs:3.6,0.12) {heavier tail};
\end{axis}
\end{tikzpicture}
\caption{Two distributions with the same mean and standard deviation. The Laplace has
more mass near the centre \emph{and} far out in the tails, with less in the shoulders.
Kurtosis is driven by the tails: a value more than 3 SDs from the mean is about five times as likely
under the Laplace ($1.4\%$ versus $0.27\%$).}
\label{fig:quantiles-kurtosis}
\end{figure}

Reference values: uniform $-1.2$, Normal $0$, Laplace $3$, and $6$ for Student's $t$
with 5 degrees of freedom (a heavy-tailed cousin of the Normal that Chapter 7 uses). Our twelve latencies have excess kurtosis $5.1$, again almost entirely from
the value 120 (it supplies over $99\%$ of $\sum (x_i - \xbar)^4$).

\begin{trapbox}
Kurtosis is often described as ``peakedness''. That is misleading. Distributions with
flat tops can have high kurtosis and pointed ones can have low kurtosis. What kurtosis
detects is how often extreme values occur. For latency data it is a warning that p99
and p99.9 will be far from the median, and that the SD (Chapter 2) is being pushed
around by a handful of points.
\end{trapbox}

\section{Taming right skew with a log transform}

Right-skewed data are hard to look at. In Figure~\ref{fig:quantiles-lat-hist} most of
the 200 requests crowd into the first few bins and the tail is a sparse scatter. Taking
logarithms fixes this. A log scale turns ratios into distances: 10 to 100 ms is the same
step as 100 to 1000 ms. Latency often builds up through multiplicative effects (a cache
miss doubles the time, a slow disk triples it), and logs turn those multiplications
into additions.

\begin{figure}[H]
\centering
\begin{tikzpicture}
\begin{axis}[
  name=raw, width=6.5cm, height=5.2cm,
  xmin=0, xmax=200, ymin=0, ymax=65, axis lines=left,
  xlabel={latency (ms)}, ylabel={requests},
  title={raw: skewness 2.44}, title style={font=\small}, axis on top,
]
\addplot[ybar interval, fill=histfill, draw=white] coordinates {
  (0,13) (10,60) (20,50) (30,25) (40,20) (50,11) (60,3) (70,7) (80,4) (90,2)
  (100,0) (110,1) (120,3) (130,0) (140,0) (150,0) (160,0) (170,0) (180,0) (190,1) (200,0)};
\end{axis}
\begin{axis}[
  at={(raw.east)}, anchor=west, xshift=1.1cm,
  width=6.5cm, height=5.2cm,
  xmin=0.4, xmax=2.4, ymin=0, ymax=32, axis lines=left,
  xlabel={$\log_{10}$ latency}, 
  title={after $\log_{10}$: skewness 0.02}, title style={font=\small}, axis on top,
  extra x ticks={1.3936}, extra x tick labels={}, extra x tick style={tick style={targetgreen, thick}},
]
\addplot[ybar interval, fill=targetgreen!45, draw=white] coordinates {
  (0.4,1) (0.5,0) (0.6,1) (0.7,2) (0.8,5) (0.9,4) (1.0,16) (1.1,21) (1.2,23) (1.3,28)
  (1.4,27) (1.5,19) (1.6,21) (1.7,13) (1.8,8) (1.9,6) (2.0,3) (2.1,1) (2.2,1) (2.3,0)};
\end{axis}
\end{tikzpicture}
\caption{The same 200 latencies before and after taking $\log_{10}$. The long right
tail becomes a roughly symmetric bell centred near $\log_{10} 25 \approx 1.4$. On the log
scale 10 ms, 100 ms and 1000 ms are equally spaced (1, 2, 3).}
\label{fig:quantiles-log}
\end{figure}

What the transform buys:
\begin{itemize}
  \item \textbf{Visibility.} The bulk and the tail fit on one readable axis.
  \item \textbf{Symmetry.} Skewness drops from $2.44$ to $0.02$, so methods that work
  best on symmetric, bell-shaped data (many in Chapters 7--13) behave better.
  \item \textbf{Interpretable effects.} A difference of $0.1$ in $\log_{10}$ latency is a
  factor of $10^{0.1} \approx 1.26$, a $26\%$ slowdown, whatever the baseline.
\end{itemize}

Back-transforming needs care. Quantiles pass straight through: the median of
$\log_{10}$ latency is $\log_{10}$ of the median, so $10^{\text{median of logs}}$ recovers
$24.8$ ms exactly. The mean of the logs, transformed back, gives the \textbf{geometric
mean}, here $25.5$ ms. It is close to the median, and well below the ordinary mean of
$32.5$ ms. You cannot recover the ordinary mean from the mean of the logs.

\begin{trapbox}
Logs need positive numbers. Zero latencies, zero counts or negative profits break
$\log x$. A common patch is $\log(x + 1)$; it works for counts but the ``$+1$'' is an
arbitrary choice that matters when values are small. Also, the log of the twelve
latencies still has skewness $1.53$: a transform cannot hide a genuine outlier such as
120 ms in a sample of 12.
\end{trapbox}

\section{QQ plots: comparing a sample with a model}

Is this sample consistent with a Normal distribution? The histogram is a blurry way
to answer, because it depends on bin width and our eye is bad at judging curves. A
\textbf{quantile--quantile plot} (QQ plot) turns the question into ``are these points
on a straight line?'', which the eye answers well.

\subsection{Building one by hand}

Eight adult heights in cm, sorted: $158, 162, 165, 167, 170, 171, 175, 184$. Mean
$169$, SD $8.07$.

The idea: if the data were a perfect Normal sample, the $i$-th smallest of $n$ values
should sit near the population quantile at fraction $(i - 0.5)/n$. (The $-0.5$ puts the
$i$-th point in the middle of its $1/n$ slice. With $n=8$ the lowest point gets
fraction $1/16$, the middle of the bottom eighth. Using $i/n$ instead would send the
top point to fraction $1$, whose Normal quantile is infinite.) So:

\begin{enumerate}
  \item Sort the data: $x_{(1)} \le \cdots \le x_{(n)}$.
  \item Compute plotting positions $p_i = (i - 0.5)/n$.
  \item Convert each to a standard Normal quantile $z_i = \Phi^{-1}(p_i)$, where $\Phi$
  is the standard Normal CDF and $\Phi^{-1}$ its quantile function.
  \item Plot the points $(z_i,\, x_{(i)})$.
\end{enumerate}

\begin{table}[H]
\centering
\begin{tabular}{cccccccccc}
\toprule
$i$ & 1 & 2 & 3 & 4 & 5 & 6 & 7 & 8 \\
\midrule
$x_{(i)}$ & 158 & 162 & 165 & 167 & 170 & 171 & 175 & 184 \\
$p_i$ & 0.0625 & 0.1875 & 0.3125 & 0.4375 & 0.5625 & 0.6875 & 0.8125 & 0.9375 \\
$z_i$ & $-1.534$ & $-0.887$ & $-0.489$ & $-0.157$ & $0.157$ & $0.489$ & $0.887$ & $1.534$ \\
\bottomrule
\end{tabular}
\caption{QQ plot ingredients for eight heights. The $z_i$ values are symmetric
because the plotting positions are.}
\label{tab:quantiles-qq}
\end{table}

If the data are Normal with mean $\mu$ and SD $\sigma$, their quantiles are
$\mu + \sigma z$, so the points should fall near the line $x = \mu + \sigma z$. Using the
sample values, that reference line is $x = 169 + 8.07z$. At $z = 1.534$ the line predicts
$169 + 8.07(1.534) \approx 181.4$; the tallest person is $184$, a bit above.

\begin{figure}[H]
\centering
\begin{tikzpicture}
\begin{axis}[
  width=9cm, height=6.6cm,
  xmin=-2, xmax=2, ymin=150, ymax=190,
  axis lines=left, xlabel={standard Normal quantile $z_i$}, ylabel={height $x_{(i)}$ (cm)},
  domain=-2:2,
]
\addplot[gray, thick, dashed] {169+8.07*x};
\addplot[only marks, mark=*, mark size=2.4pt, sourceblue] coordinates {
  (-1.534,158) (-0.887,162) (-0.489,165) (-0.157,167) (0.157,170) (0.489,171) (0.887,175) (1.534,184)};
\draw[densityred, thick] (axis cs:1.534,181.4) -- (axis cs:1.534,184);
\node[densityred, anchor=west, font=\scriptsize] at (axis cs:1.6,182.7) {$+2.6$};
\node[gray, anchor=west, font=\small] at (axis cs:0.3,156) {line $169 + 8.07z$};
\end{axis}
\end{tikzpicture}
\caption{QQ plot of eight heights against the Normal. The points hug the reference line,
so a Normal model is reasonable. With $n = 8$ wobbles of a few cm are expected.}
\label{fig:quantiles-qq-hand}
\end{figure}

The slope of the line is the SD and its height at $z=0$ is the mean. That is why the
QQ plot compares shape only: a Normal sample with any mean and SD gives a straight
line, and only departures from straightness signal a different shape.

\subsection{Reading the patterns}

Departures from the line are easiest to read at the ends, where the tails live.
Figure~\ref{fig:quantiles-qq-patterns} shows the idealised pattern for four shapes,
each standardised to mean 0 and SD 1 and plotted with $n=40$ points.

\begin{figure}[H]
\centering
\begin{tikzpicture}
\pgfplotsset{qqpanel/.style={width=6.2cm, height=4cm, xmin=-2.6, xmax=2.6, ymin=-3.2, ymax=3.8,
  axis lines=left, domain=-2.6:2.6, title style={font=\small}, tick label style={font=\scriptsize}}}
\begin{axis}[qqpanel, name=a, title={Normal: straight}]
\addplot[gray, dashed] {x};
\addplot[only marks, mark=*, mark size=1.2pt, sourceblue] coordinates {(-2.24,-2.28) (-1.78,-1.81) (-1.53,-1.56) (-1.36,-1.38) (-1.21,-1.23) (-1.09,-1.11) (-0.98,-1.0) (-0.89,-0.9) (-0.8,-0.81) (-0.71,-0.73) (-0.64,-0.65) (-0.56,-0.57) (-0.49,-0.5) (-0.42,-0.43) (-0.35,-0.36) (-0.29,-0.29) (-0.22,-0.22) (-0.16,-0.16) (-0.09,-0.1) (-0.03,-0.03) (0.03,0.03) (0.09,0.1) (0.16,0.16) (0.22,0.22) (0.29,0.29) (0.35,0.36) (0.42,0.43) (0.49,0.5) (0.56,0.57) (0.64,0.65) (0.71,0.73) (0.8,0.81) (0.89,0.9) (0.98,1.0) (1.09,1.11) (1.21,1.23) (1.36,1.38) (1.53,1.56) (1.78,1.81) (2.24,2.28)};
\end{axis}
\begin{axis}[qqpanel, name=b, at={(a.east)}, anchor=west, xshift=0.6cm, title={heavy tails ($t_3$): S-shape}]
\addplot[gray, dashed] {x};
\addplot[only marks, mark=*, mark size=1.2pt, densityred] coordinates {(-2.24,-2.88) (-1.78,-1.85) (-1.53,-1.46) (-1.36,-1.22) (-1.21,-1.05) (-1.09,-0.92) (-0.98,-0.81) (-0.89,-0.72) (-0.8,-0.64) (-0.71,-0.56) (-0.64,-0.49) (-0.56,-0.43) (-0.49,-0.37) (-0.42,-0.32) (-0.35,-0.27) (-0.29,-0.22) (-0.22,-0.17) (-0.16,-0.12) (-0.09,-0.07) (-0.03,-0.02) (0.03,0.02) (0.09,0.07) (0.16,0.12) (0.22,0.17) (0.29,0.22) (0.35,0.27) (0.42,0.32) (0.49,0.37) (0.56,0.43) (0.64,0.49) (0.71,0.56) (0.8,0.64) (0.89,0.72) (0.98,0.81) (1.09,0.92) (1.21,1.05) (1.36,1.22) (1.53,1.46) (1.78,1.85) (2.24,2.88)};
\end{axis}
\begin{axis}[qqpanel, name=c, at={(a.south)}, anchor=north, yshift=-1.5cm, title={right skew: bends up}]
\addplot[gray, dashed] {x};
\addplot[only marks, mark=*, mark size=1.2pt, quotientorange] coordinates {(-2.24,-1.02) (-1.78,-1.0) (-1.53,-0.97) (-1.36,-0.94) (-1.21,-0.91) (-1.09,-0.88) (-0.98,-0.85) (-0.89,-0.82) (-0.8,-0.79) (-0.71,-0.75) (-0.64,-0.72) (-0.56,-0.68) (-0.49,-0.64) (-0.42,-0.61) (-0.35,-0.57) (-0.29,-0.52) (-0.22,-0.48) (-0.16,-0.43) (-0.09,-0.39) (-0.03,-0.34) (0.03,-0.29) (0.09,-0.23) (0.16,-0.17) (0.22,-0.11) (0.29,-0.05) (0.35,0.02) (0.42,0.1) (0.49,0.18) (0.56,0.27) (0.64,0.36) (0.71,0.47) (0.8,0.58) (0.89,0.71) (0.98,0.86) (1.09,1.04) (1.21,1.25) (1.36,1.51) (1.53,1.86) (1.78,2.4) (2.24,3.54)};
\end{axis}
\begin{axis}[qqpanel, at={(c.east)}, anchor=west, xshift=0.6cm, title={light tails (uniform): flat}]
\addplot[gray, dashed] {x};
\addplot[only marks, mark=*, mark size=1.2pt, structurepurple] coordinates {(-2.24,-1.69) (-1.78,-1.6) (-1.53,-1.52) (-1.36,-1.43) (-1.21,-1.34) (-1.09,-1.26) (-0.98,-1.17) (-0.89,-1.08) (-0.8,-1.0) (-0.71,-0.91) (-0.64,-0.82) (-0.56,-0.74) (-0.49,-0.65) (-0.42,-0.56) (-0.35,-0.48) (-0.29,-0.39) (-0.22,-0.3) (-0.16,-0.22) (-0.09,-0.13) (-0.03,-0.04) (0.03,0.04) (0.09,0.13) (0.16,0.22) (0.22,0.3) (0.29,0.39) (0.35,0.48) (0.42,0.56) (0.49,0.65) (0.56,0.74) (0.64,0.82) (0.71,0.91) (0.8,1.0) (0.89,1.08) (0.98,1.17) (1.09,1.26) (1.21,1.34) (1.36,1.43) (1.53,1.52) (1.78,1.6) (2.24,1.69)};
\end{axis}
\end{tikzpicture}
\caption{Four QQ plot signatures against the Normal (horizontal axis: Normal quantile;
vertical: standardised data). Points are the exact quantiles of each distribution, so
the patterns are noise-free; real samples wobble around them, most at the ends.}
\label{fig:quantiles-qq-patterns}
\end{figure}

\begin{table}[H]
\centering
\begin{tabular}{ll}
\toprule
What you see & What it means \\
\midrule
Points on the line & Shape matches the Normal \\
Low end below, high end above the line (S) & Heavier tails than Normal: more extreme values \\
Low end above, high end below the line & Lighter tails than Normal: values bounded \\
Curve bending upward (convex) & Right skew: long upper tail \\
Curve bending downward (concave) & Left skew: long lower tail \\
One or two points far off, rest on the line & Isolated outliers \\
\bottomrule
\end{tabular}
\caption{Reading a Normal QQ plot.}
\label{tab:quantiles-qq-read}
\end{table}

The logic behind every row is the same. The vertical position of a point is the data
quantile; the line marks where a Normal quantile would be. In a heavy-tailed sample the
extreme quantiles are more extreme than Normal ones, so the ends of the plot pull away
from the line, downward at the left and upward at the right. In a right-skewed sample the
low quantiles are bunched (above the line) and the high quantiles spread out (above the
line again, increasingly), giving a curve that bends up.

\begin{lstlisting}[language=Python]
import numpy as np
from scipy import stats
x = np.sort(np.array([158,162,165,167,170,171,175,184]))
p = (np.arange(1, len(x) + 1) - 0.5) / len(x)
z = stats.norm.ppf(p)                # -1.534 ... 1.534
# plot (z, x); or let scipy do it all:
stats.probplot(x, dist="norm")       # uses slightly different p_i
\end{lstlisting}

\begin{connectionbox}
QQ plots are not only for the Normal. Replace $\Phi^{-1}$ with any distribution's
quantile function to test that model; plot log-latency against the Normal to check a
log-normal model. Once regression arrives (Chapter 12), a Normal QQ plot of the
residuals is the standard check of the Normal-errors assumption.
\end{connectionbox}

\begin{recoverybox}
A QQ plot of request latencies against the Normal curves sharply upward at the right
end. You take logs and redraw; now the points lie on a line. What distribution do the
raw latencies follow, approximately? (Log-normal: their logarithm is Normal.)
\end{recoverybox}

\section{Smooth density estimates}

A histogram's shape depends on where you put the bins. Shift the bin edges by half a
bin and a bump can appear or vanish. A \textbf{kernel density estimate} (KDE) replaces
the boxes with smooth bumps and removes the bin-edge arbitrariness.

\subsection{Building a KDE from bumps}

Place a small bell curve, called a \textbf{kernel}, centred on each data point. Every
bell has the same width $h$, called the \textbf{bandwidth}, and area $1/n$. Add the
bells up. Where data points crowd together their bells pile up into a hill; where data
are sparse, the sum is low.

\begin{definition}[Kernel density estimate]
With a kernel $K$ (a symmetric density, usually the standard Normal $\varphi$) and
bandwidth $h > 0$,
\[
  \hat f_h(x) = \frac{1}{nh}\sum_{i=1}^n K\!\left(\frac{x - x_i}{h}\right).
\]
\end{definition}

The factor $1/h$ inside and outside keeps each bump's area at $1/n$ as you stretch it,
so the total area is $1$: $\hat f_h$ is a genuine density. At a point $x$, the estimate
counts nearby data, with closer points counting more.

\begin{figure}[H]
\centering
\begin{tikzpicture}
\begin{axis}[
  width=12cm, height=5.6cm,
  xmin=-1, xmax=14, ymin=0, ymax=0.2,
  axis lines=left, xlabel={$x$}, ylabel={density},
  domain=-1:14, samples=151,
  legend style={font=\small}, legend pos=north west,
  scaled y ticks=false, yticklabel style={/pgf/number format/fixed},
]
\addplot[densityred, very thick] {(exp(-(x-2)^2/2)+exp(-(x-3.5)^2/2)+exp(-(x-4)^2/2)+exp(-(x-7)^2/2)+exp(-(x-7.5)^2/2)+exp(-(x-8)^2/2)+exp(-(x-11)^2/2))*0.39894/7};
\addlegendentry{KDE (sum of bumps)}
\addplot[distblue, thin] {0.39894/7*exp(-(x-2)^2/2)};
\addlegendentry{one kernel, area $1/7$}
\foreach \c in {3.5,4,7,7.5,8,11} {
  \addplot[distblue, thin] {0.39894/7*exp(-(x-\c)^2/2)};
}
\addplot[only marks, mark=|, mark size=4pt, thick, black] coordinates {(2,0) (3.5,0) (4,0) (7,0) (7.5,0) (8,0) (11,0)};
\end{axis}
\end{tikzpicture}
\caption{A KDE of seven points with Gaussian kernels of bandwidth $h=1$. Each thin bump
has area $1/7$. The thick curve is their sum. The three points near 7.5 build the
tallest hill.}
\label{fig:quantiles-kde-build}
\end{figure}

At $x = 7$ in Figure~\ref{fig:quantiles-kde-build}, the three bumps at $7$, $7.5$ and $8$
contribute most; the estimate there is
$\hat f(7) = \frac17\sum_i \varphi(7 - x_i) \approx 0.143$.

\subsection{Bandwidth is the whole game}

The kernel shape barely matters. The bandwidth matters enormously. Small $h$ gives narrow
bumps and a spiky curve that follows every quirk of the sample. Large $h$ gives wide
bumps that blur real features away. Figure~\ref{fig:quantiles-kde-bw} shows 38
reaction times drawn from two groups (a fast majority near 30 ms and a slower group near
50 ms).

\begin{figure}[H]
\centering
\begin{tikzpicture}
\begin{axis}[
  width=12.5cm, height=6cm,
  xmin=10, xmax=75, ymin=0, ymax=0.105,
  axis lines=left, xlabel={reaction time (ms)}, ylabel={density},
  legend style={font=\small}, legend pos=north east,
  scaled y ticks=false, yticklabel style={/pgf/number format/fixed},
]
\addplot[gray, thin] coordinates {(10.0,0.0000) (11.0,0.0000) (12.0,0.0000) (13.0,0.0000) (14.0,0.0000) (15.0,0.0000) (16.0,0.0000) (17.0,0.0000) (18.0,0.0000) (19.0,0.0000) (20.0,0.0002) (21.0,0.0019) (22.0,0.0090) (23.0,0.0211) (24.0,0.0291) (25.0,0.0274) (26.0,0.0239) (27.0,0.0239) (28.0,0.0293) (29.0,0.0541) (30.0,0.0834) (31.0,0.0838) (32.0,0.0628) (33.0,0.0451) (34.0,0.0418) (35.0,0.0388) (36.0,0.0287) (37.0,0.0186) (38.0,0.0131) (39.0,0.0121) (40.0,0.0079) (41.0,0.0036) (42.0,0.0076) (43.0,0.0168) (44.0,0.0224) (45.0,0.0197) (46.0,0.0115) (47.0,0.0139) (48.0,0.0299) (49.0,0.0426) (50.0,0.0458) (51.0,0.0423) (52.0,0.0262) (53.0,0.0100) (54.0,0.0082) (55.0,0.0106) (56.0,0.0072) (57.0,0.0075) (58.0,0.0106) (59.0,0.0064) (60.0,0.0014) (61.0,0.0001) (62.0,0.0000) (63.0,0.0000) (64.0,0.0000) (65.0,0.0000) (66.0,0.0000) (67.0,0.0000) (68.0,0.0000) (69.0,0.0000) (70.0,0.0000) (71.0,0.0000) (72.0,0.0000) (73.0,0.0000) (74.0,0.0000) (75.0,0.0000)};
\addlegendentry{$h=1$: too wiggly}
\addplot[densityred, very thick] coordinates {(10.0,0.0000) (11.0,0.0001) (12.0,0.0001) (13.0,0.0003) (14.0,0.0005) (15.0,0.0009) (16.0,0.0015) (17.0,0.0025) (18.0,0.0039) (19.0,0.0057) (20.0,0.0081) (21.0,0.0110) (22.0,0.0144) (23.0,0.0183) (24.0,0.0226) (25.0,0.0271) (26.0,0.0318) (27.0,0.0363) (28.0,0.0404) (29.0,0.0437) (30.0,0.0458) (31.0,0.0466) (32.0,0.0458) (33.0,0.0436) (34.0,0.0402) (35.0,0.0361) (36.0,0.0316) (37.0,0.0273) (38.0,0.0235) (39.0,0.0204) (40.0,0.0182) (41.0,0.0170) (42.0,0.0166) (43.0,0.0171) (44.0,0.0182) (45.0,0.0197) (46.0,0.0213) (47.0,0.0229) (48.0,0.0240) (49.0,0.0246) (50.0,0.0244) (51.0,0.0234) (52.0,0.0217) (53.0,0.0195) (54.0,0.0170) (55.0,0.0144) (56.0,0.0119) (57.0,0.0097) (58.0,0.0077) (59.0,0.0060) (60.0,0.0046) (61.0,0.0034) (62.0,0.0024) (63.0,0.0017) (64.0,0.0011) (65.0,0.0007) (66.0,0.0004) (67.0,0.0002) (68.0,0.0001) (69.0,0.0001) (70.0,0.0000) (71.0,0.0000) (72.0,0.0000) (73.0,0.0000) (74.0,0.0000) (75.0,0.0000)};
\addlegendentry{$h=4$: two modes visible}
\addplot[distblue, thick, dashed] coordinates {(10.0,0.0055) (11.0,0.0062) (12.0,0.0070) (13.0,0.0079) (14.0,0.0088) (15.0,0.0098) (16.0,0.0108) (17.0,0.0119) (18.0,0.0129) (19.0,0.0140) (20.0,0.0151) (21.0,0.0162) (22.0,0.0173) (23.0,0.0184) (24.0,0.0194) (25.0,0.0203) (26.0,0.0212) (27.0,0.0220) (28.0,0.0228) (29.0,0.0234) (30.0,0.0240) (31.0,0.0244) (32.0,0.0248) (33.0,0.0250) (34.0,0.0252) (35.0,0.0252) (36.0,0.0251) (37.0,0.0250) (38.0,0.0247) (39.0,0.0244) (40.0,0.0240) (41.0,0.0236) (42.0,0.0230) (43.0,0.0225) (44.0,0.0218) (45.0,0.0211) (46.0,0.0204) (47.0,0.0197) (48.0,0.0189) (49.0,0.0181) (50.0,0.0173) (51.0,0.0165) (52.0,0.0156) (53.0,0.0148) (54.0,0.0140) (55.0,0.0131) (56.0,0.0123) (57.0,0.0115) (58.0,0.0107) (59.0,0.0099) (60.0,0.0091) (61.0,0.0083) (62.0,0.0076) (63.0,0.0069) (64.0,0.0063) (65.0,0.0056) (66.0,0.0050) (67.0,0.0045) (68.0,0.0040) (69.0,0.0035) (70.0,0.0031) (71.0,0.0027) (72.0,0.0023) (73.0,0.0020) (74.0,0.0017) (75.0,0.0015)};
\addlegendentry{$h=12$: second mode gone}
\addplot[only marks, mark=|, mark size=3pt, black] coordinates {(22.9,0) (24.0,0) (24.4,0) (25.6,0) (26.8,0) (27.8,0) (29.4,0) (29.5,0) (29.8,0) (29.9,0) (30.1,0) (30.4,0) (30.7,0) (31.2,0) (31.3,0) (31.5,0) (31.9,0) (32.5,0) (33.4,0) (34.0,0) (34.5,0) (35.2,0) (35.7,0) (37.1,0) (39.2,0) (43.0,0) (44.2,0) (45.1,0) (47.9,0) (48.5,0) (49.2,0) (49.4,0) (50.3,0) (50.9,0) (51.1,0) (51.8,0) (54.9,0) (58.0,0)};
\end{axis}
\end{tikzpicture}
\caption{Three KDEs of the same 38 points (tick marks at the bottom). $h=1$ invents
bumps that are just sampling noise. $h=12$ smooths away the slower group. $h=4$, a little
below Silverman's rule of thumb ($5.1$), shows the two groups.}
\label{fig:quantiles-kde-bw}
\end{figure}

This is the first appearance of a trade-off that runs through all of statistics and
machine learning. A small bandwidth has low \emph{bias} (it can follow any shape) but
high \emph{variance} (a new sample would give a very different curve). A large bandwidth
is stable but systematically wrong. Software picks $h$ with a rule of thumb; the most
common, \textbf{Silverman's rule}, is
\[
  h = 1.06\, s\, n^{-1/5},
\]
where $s$ is the sample SD. For these 38 points ($s \approx 9.9$) it gives $h \approx 1.06 \times 9.9 \times 38^{-1/5}
\approx 5.1$.
The rule is tuned for roughly Normal data and tends to oversmooth data with several
modes, so always try a smaller and a larger value and see which features survive.

\begin{lstlisting}[language=Python]
from scipy import stats
kde = stats.gaussian_kde(x)              # Scott's rule, close to Silverman's
kde_narrow = stats.gaussian_kde(x, bw_method=0.1)   # factor times SD of x
grid = np.linspace(10, 75, 200)
density = kde(grid)                      # evaluate the curve on a grid
\end{lstlisting}

\begin{trapbox}
A KDE puts probability where no data can be. Gaussian bumps around latencies near zero
spill into negative milliseconds; bumps around exam scores of 100 spill past 100. For
bounded data, either estimate on a transformed scale (log for positive data) or read the
curve near the boundary with suspicion.
\end{trapbox}

\subsection{Violin plots}

A \textbf{violin plot} is a box plot with a KDE wrapped around it: the density is drawn
vertically and mirrored, so wide parts of the violin are where data are dense. It keeps
the box plot's compactness for comparing groups and adds the shape information the box
plot throws away.

\begin{figure}[H]
\centering
\begin{tikzpicture}
\begin{axis}[
  width=8cm, height=6.6cm,
  xmin=0.4, xmax=2.6, ymin=15, ymax=65,
  xtick={1,2}, xticklabels={box plot, violin plot},
  ylabel={reaction time (ms)}, axis lines=left,
  boxplot/draw direction=y,
]
\addplot[fill=targetgreen!25, draw=targetgreen!70!black, thick] coordinates {
  (2.022,17.0) (2.041,18.5) (2.069,20.0) (2.108,21.5) (2.156,22.9) (2.211,24.4) (2.271,25.9) (2.328,27.4) (2.375,28.9) (2.400,30.4) (2.398,31.8) (2.368,33.3) (2.319,34.8) (2.262,36.3) (2.210,37.8) (2.171,39.3) (2.149,40.7) (2.144,42.2) (2.154,43.7) (2.173,45.2) (2.194,46.7) (2.209,48.2) (2.212,49.6) (2.201,51.1) (2.176,52.6) (2.145,54.1) (2.112,55.6) (2.082,57.1) (2.058,58.5) (2.039,60.0) (2.025,61.5) (2.015,63.0) (1.985,63.0) (1.975,61.5) (1.961,60.0) (1.942,58.5) (1.918,57.1) (1.888,55.6) (1.855,54.1) (1.824,52.6) (1.799,51.1) (1.788,49.6) (1.791,48.2) (1.806,46.7) (1.827,45.2) (1.846,43.7) (1.856,42.2) (1.851,40.7) (1.829,39.3) (1.790,37.8) (1.738,36.3) (1.681,34.8) (1.632,33.3) (1.602,31.8) (1.600,30.4) (1.625,28.9) (1.672,27.4) (1.729,25.9) (1.789,24.4) (1.844,22.9) (1.892,21.5) (1.931,20.0) (1.959,18.5) (1.978,17.0)} -- cycle;
\addplot+[boxplot prepared={lower whisker=22.9, lower quartile=29.95, median=33.7,
  upper quartile=47.2, upper whisker=58.0, draw position=1, box extend=0.35},
  draw=sourceblue, fill=histfill!50, thick, mark=none] coordinates {};
\addplot+[boxplot prepared={lower whisker=22.9, lower quartile=29.95, median=33.7,
  upper quartile=47.2, upper whisker=58.0, draw position=2, box extend=0.08},
  draw=black, fill=white, thick, mark=none] coordinates {};
\end{axis}
\end{tikzpicture}
\caption{The 38 reaction times as a box plot and as a violin plot. The box plot shows a
tall upper half but cannot show that the data form two separate groups. The violin's
two bulges make the two groups obvious.}
\label{fig:quantiles-violin}
\end{figure}

\begin{historybox}
John Tukey (1915--2000), a statistician at Princeton and Bell Labs, spent the 1960s and
1970s arguing that statisticians jumped to formal models before looking at their data.
His problem was practical: with pencil and paper, how do you summarise a batch of
numbers in a way that a few wild values cannot distort? His answer, in the 1977 book
\emph{Exploratory Data Analysis}, used medians and quartiles (which ignore how extreme
the extremes are) instead of means and SDs, and drew them as the box-and-whisker plot,
with the $1.5\,\IQR$ fences to flag points worth a second look. It unlocked fast visual
comparison of many groups at once, and it is still the default summary plot in every
statistics package. Violin plots (Hintze and Nelson, 1998) and the ECDF-based tools in
this chapter are refinements of the same idea: look at the whole distribution before
summarising it.
\end{historybox}

\section{Choosing a picture}

Each display in this chapter keeps different information.

\begin{figure}[H]
\centering
\begin{tikzpicture}[
  node distance=0.55cm and 0.5cm,
  q/.style={draw, rounded corners, fill=quotientorange!15, align=center, font=\small, text width=3.4cm},
  a/.style={draw, rounded corners, fill=sourceblue!12, align=center, font=\small, text width=3.3cm},
  >={Stealth}
]
\node[q] (start) {What do you need to see?};
\node[a, below left=0.7cm and 1.6cm of start] (box) {Box plot or violin\\(compare groups)};
\node[a, below=0.7cm of start] (ecdf) {ECDF\\(percentiles, SLOs)};
\node[a, below right=0.7cm and 1.6cm of start] (qq) {QQ plot\\(fit to a model)};
\node[a, below=0.6cm of ecdf] (kde) {Histogram or KDE\\(shape, modes)};
\node[a, below=0.6cm of box] (log) {Skewed? Plot on log scale first};
\draw[->] (start) -- node[above left, font=\scriptsize, pos=0.5] {many groups} (box);
\draw[->] (start) -- node[right, font=\scriptsize] {thresholds} (ecdf);
\draw[->] (start) -- node[above right, font=\scriptsize, pos=0.5] {is it Normal?} (qq);
\draw[->] (ecdf) -- node[right, font=\scriptsize] {bumps, gaps} (kde);
\draw[->, dashed] (box) -- (log);
\end{tikzpicture}
\caption{A rough guide to choosing a display. In practice look at two: one that shows
the whole shape (ECDF, KDE) and one that summarises (box plot, a few percentiles).}
\label{fig:quantiles-choose}
\end{figure}

\section{Summary}

\begin{itemize}
  \item The $p$-quantile has a fraction $p$ of the data below it. With finite data the
  definition needs an interpolation rule; numpy offers nine. They agree in the middle and
  can disagree badly in sparse tails.
  \item Five-number summary: min, $Q_1$, median, $Q_3$, max. $\IQR = Q_3 - Q_1$.
  \item Box plot: box from $Q_1$ to $Q_3$, median line, whiskers to the last points within
  $1.5\,\IQR$ of the box, outliers plotted individually. Roughly $0.7\%$ of Normal data and
  $4.8\%$ of exponential data fall outside the fences.
  \item ECDF $\hat F(t)$ = fraction of data $\le t$. Quantiles are read by going across
  and down. The quantile function is the inverse of the CDF.
  \item Quantiles commute with increasing transforms: $Q_{g(X)}(p) = g(Q_X(p))$. Means do not.
  \item For latency, report p50/p90/p99 with counts. The mean hides tails. Percentiles do
  not average across servers or time windows; merge histograms or raw data instead.
  \item Skewness $g_1 = m_3/m_2^{3/2}$ measures tail imbalance; excess kurtosis
  $m_4/m_2^2 - 3$ measures tail weight relative to the Normal. Both are sensitive to single
  points.
  \item Logs turn right-skewed, multiplicative data into roughly symmetric data. The
  back-transformed mean of logs is the geometric mean.
  \item QQ plot: sorted data against $\Phi^{-1}((i-0.5)/n)$. Straight means Normal-shaped;
  S-shape means heavy tails; upward bend means right skew.
  \item KDE: a sum of bumps of width $h$. Bandwidth controls the bias--variance trade-off.
  Violin plots wrap a KDE around a box plot.
\end{itemize}

Chapter 4 moves from one variable to two: scatter plots, covariance and correlation.

\section{Exercises}

\begin{enumerate}
  \item For the data $3, 7, 8, 12, 20$, compute $Q_1$, the median and $Q_3$ with numpy's
  default (linear) method. Then compute the fences and decide whether any point is an
  outlier.
  \item For the same data, compute $Q_1$ with the \texttt{weibull} rule (position
  $i/(n+1)$).
  \item Using the twelve latencies of this chapter, find $\hat F(25)$ and $\hat F(40)$.
  \item Latency is exponential with mean 50 ms. Find the median and the p99.
  \item A page makes 20 independent backend calls and waits for all of them. What is the
  probability that at least one is slower than the backend's p99?
  \item Server A served requests of $10, 20, 30$ ms; server B served one request of
  $100$ ms. Compare the average of the two medians with the median of all four requests.
  \item The median latency of a service is 25 ms. What is the median of $\log_{10}$
  latency? Can you say what the mean of $\log_{10}$ latency is?
  \item A Normal QQ plot shows points above the line at the left end and below the line at
  the right end. What does that say about the tails? What shape would a curve that bends
  downward (concave) indicate?
  \item Out of 1000 values drawn from a Normal distribution, about how many would a box plot
  flag as outliers?
  \item Describe $\hat f_h$ in the limits $h \to 0$ and $h \to \infty$.
\end{enumerate}

\subsection*{Answers}

\begin{enumerate}
  \item $n=5$, $h = 4p + 1$. $Q_1$: $h=2$, so $Q_1 = 7$. Median $8$. $Q_3$: $h = 4$, so
  $Q_3 = 12$. $\IQR = 5$, fences $7 - 7.5 = -0.5$ and $12 + 7.5 = 19.5$. The value $20$ is
  an outlier.
  \item $h = (n+1)p = 6(0.25) = 1.5$, so $Q_1 = 3 + 0.5(7 - 3) = 5$.
  \item $\hat F(25) = 9/12 = 0.75$; $\hat F(40) = 10/12 \approx 0.83$.
  \item Median $= 50\ln 2 \approx 34.7$ ms; p99 $= 50\ln 100 \approx 230.3$ ms.
  \item $1 - 0.99^{20} \approx 0.18$.
  \item Medians 20 and 100 average to 60. The four requests $10, 20, 30, 100$ have median
  $25$.
  \item $\log_{10} 25 \approx 1.40$, because quantiles pass through increasing transforms.
  The mean of the logs cannot be found from the median alone (it is $\log_{10}$ of the
  geometric mean).
  \item Light tails (the extremes are less extreme than Normal ones), as in uniform data.
  A concave curve indicates left skew.
  \item About $0.7\%$ of 1000, so roughly 7.
  \item As $h \to 0$ the bumps become spikes at the data points (the curve is zero almost
  everywhere and infinitely tall at each $x_i$). As $h \to \infty$ the bumps become one
  very wide, flat curve that forgets the data's shape.
\end{enumerate}
